\documentclass[a4paper,11pt]{article}

\usepackage[utf8]{inputenc}
\usepackage[T1]{fontenc}
\usepackage[french]{babel}
\usepackage{amsmath,amssymb}
\usepackage{hyperref}
\usepackage{geometry}
\geometry{margin=2.5cm}

\title{Rapport sur le mini-projet MapReduce\newline Analyse des critiques Amazon Movies \\ (Notebook ``Moubarak\_LASSISSI\_\_02\_Exe\_MapReduce'')}
\author{Moubarak LASSISSI}
\date{\today}

\begin{document}

\maketitle

\tableofcontents

\section{Introduction}

Ce mini-projet a pour objectif de mettre en pratique le paradigme
\textbf{MapReduce} en utilisant un jeu de données réel~: les critiques
(Amazon Reviews) de la catégorie \emph{Movies \\& TV}. Dans le rôle
d'un data engineer pour une plateforme de streaming, le but est de~:
\begin{itemize}
  \item compter le nombre total de critiques par film,\
  \item calculer la note moyenne par film,\
  \item extraire des mots-clés fréquents dans les textes de critiques,\
  \item joindre les notes moyennes avec les mots-clés dominants par film.
\end{itemize}

Tout cela est réalisé avec des programmes \texttt{mapper} / \texttt{reducer}
écrits en Python (mode \emph{Hadoop Streaming}) et exécutés en local
sur un environnement type Colab/Hadoop.

Ce rapport présente d'abord un rappel du modèle MapReduce, puis décrit
le jeu de données utilisé et les différentes tâches codées dans le
notebook.

\section{Rappel sur le modèle MapReduce}

\subsection{Principe général}

MapReduce est un modèle de programmation distribué, introduit par Google,
qui permet de traiter de gros volumes de données de façon parallèle
sur un cluster.

Un job MapReduce comporte classiquement deux phases principales~:
\begin{itemize}
  \item \textbf{Map}~: applique une fonction à chaque enregistrement d'entrée
        et émet des paires clé--valeur \((k, v)\).\
  \item \textbf{Reduce}~: agrège les valeurs associées à une même clé~: pour
        une clé \(k\), la fonction de reduce reçoit la liste de toutes les
        valeurs associées et produit un ou plusieurs résultats agrégés.
\end{itemize}

Entre ces deux phases, le framework Hadoop se charge automatiquement
du \emph{shuffle} et du \emph{sort}~: les paires émises par les mappers
sont triées par clé et réparties vers les reducers correspondants.

\subsection{Hadoop Streaming et scripts Python}

Dans ce projet, les mappers et reducers sont écrits en Python et
utilisés via \textbf{Hadoop Streaming}. Dans ce mode~:
\begin{itemize}
  \item le mapper lit les données depuis l'entrée standard (\texttt{stdin})
        ligne par ligne et écrit ses sorties sur la sortie standard
        (\texttt{stdout}),\
  \item le reducer lit les sorties triées des mappers sur \texttt{stdin},
        agrège par clé et écrit les résultats finaux sur \texttt{stdout}.
\end{itemize}

Cette approche permet de bénéficier de l'orchestration Hadoop en restant
sur un langage familier (Python).

\section{Jeu de données et préparation de l'environnement}

\subsection{Jeu de données Amazon Movies \\& TV 5-core}

Le notebook utilise le jeu de données \emph{Amazon Movies \\& TV 5-core},
disponible publiquement. Chaque enregistrement est un objet JSON avec
notamment les champs~:
\begin{itemize}
  \item \texttt{reviewerID}~: identifiant du critique,\
  \item \texttt{asin}~: identifiant unique du produit (ici un film),\
  \item \texttt{reviewText}~: texte complet de la critique,\
  \item \texttt{overall}~: note globale (entre 1 et 5 étoiles),\
  \item \texttt{vote}~: nombre de votes utiles pour la critique,\
  \item \texttt{category}~: ici toujours ``Movies \\& TV''.
\end{itemize}

Les données brutes sont stockées dans un fichier compressé~:\\
\texttt{Movies\_and\_TV\_5.json.gz} où chaque ligne est un objet JSON.

\subsection{Installation locale de Hadoop}

Dans les premières cellules, l'environnement est préparé pour
exécuter Hadoop en mode local~:
\begin{itemize}
  \item installation de Java 8 (\texttt{openjdk-8-jdk-headless}),\
  \item téléchargement de Hadoop~3.3.6 depuis le site officiel,\
  \item décompression de l'archive,\
  \item définition des variables d'environnement \texttt{JAVA\_HOME},
        \texttt{HADOOP\_HOME} et mise à jour du \texttt{PATH}.
\end{itemize}

Cela permet ensuite d'utiliser les commandes Hadoop (par exemple
\texttt{hadoop jar ...} ou le binaire de streaming) si nécessaire.

\subsection{Chargement et conversion des données}

Le fichier compressé \texttt{Movies\_and\_TV\_5.json.gz} est téléchargé
puis lu en Python avec le module \texttt{gzip}. Le code~:
\begin{itemize}
  \item parcourt chaque ligne,\
  \item effectue un \texttt{json.loads(line)} pour transformer la ligne
        en dictionnaire Python,\
  \item gère les erreurs de parsing avec un \texttt{try/except} pour
        ignorer les lignes malformées,\
  \item accumule les enregistrements dans une liste.
\end{itemize}

Ensuite, afin de faciliter l'utilisation en mode streaming, les données
sont réécrites dans un fichier \texttt{movies.json} au format
\textbf{JSON Lines} (une ligne = un objet JSON), en limitant à
900~000 enregistrements pour réduire les temps d'exécution en environnement
Colab. Sur un vrai cluster, cette limitation n'est pas nécessaire.

Enfin, quelques lignes de \texttt{movies.json} sont affichées pour
vérifier la structure des enregistrements.

\section{Tâche 1~: nombre total de critiques par film}

\subsection{Objectif}

La première tâche vise à compter le nombre total de critiques
(reviews) pour chaque film identifié par son code \texttt{asin}.
Il s'agit d'un équivalent \emph{movie count} d'un classique
\emph{word count}.

\subsection{Mapper de la tâche 1}

Le fichier \texttt{mapper\_task1.py} contient la fonction \texttt{mapper()}~:
\begin{itemize}
  \item lecture ligne par ligne sur \texttt{stdin},\
  \item pour chaque ligne non vide, tentative de \texttt{json.loads},\
  \item extraction de \texttt{asin},\
  \item si l'asin est présent, émission de la paire
        \texttt{"asin\textbackslash t1"}.
\end{itemize}

Autrement dit, pour chaque critique, le mapper émet
\((\text{asin}, 1)\).

\subsection{Reducer de la tâche 1}

Le fichier \texttt{reduce\_task1.py} implémente la fonction
\texttt{reducer()}~:
\begin{itemize}
  \item les lignes sont supposées triées par clé (\texttt{asin}),\
  \item on maintient une variable \texttt{current\_asin} et un compteur
        \texttt{current\_count},\
  \item à chaque changement d'\texttt{asin}, on affiche
        \texttt{"current\_asin\textbackslash tcurrent\_count"} et on
        réinitialise le compteur pour la nouvelle clé,\
  \item à la fin, on affiche le résultat pour la dernière clé.
\end{itemize}

Le résultat final pour chaque film est donc~:
\[\text{asin} \quad \text{total\_reviews}.\]

\subsection{Tests et tri des résultats}

Pour tester la chaîne mapper--reducer, le notebook utilise~:
\begin{itemize}
  \item un échantillon de 1000 lignes de \texttt{movies.json}
        (avec \texttt{head -1000}),\
  \item un tri intermédiaire avec la commande Unix \texttt{sort},\
  \item l'exécution du reducer et l'affichage de quelques lignes
        de sortie (\texttt{head -5}).
\end{itemize}

Ensuite, la chaîne complète est appliquée à l'ensemble du fichier
pour générer \texttt{task1\_output.txt}, qui est ensuite trié
par nombre de critiques décroissant afin de lister les films
les plus critiqués.

\section{Tâche 2~: note moyenne par film}

\subsection{Objectif}

La deuxième tâche calcule la \textbf{note moyenne} pour chaque film
à partir du champ \texttt{overall}. Pour un film donné, on veut~:
\[
  \text{average\_rating} = \frac{\sum_{i=1}^{n} \text{rating}_i}{n}
\]

\subsection{Mapper de la tâche 2}

Le fichier \texttt{mapper\_task2.py} lit chaque ligne JSON, en extrait~:
\begin{itemize}
  \item \texttt{asin},\
  \item \texttt{overall} (la note).
\end{itemize}

Si ces deux champs sont présents, le mapper émet la paire
\texttt{"asin\textbackslash toverall"}.

\subsection{Reducer de la tâche 2}

Le reducer \texttt{reduce\_task2.py} agrège les notes par film~:
\begin{itemize}
  \item comme pour la tâche 1, les entrées sont triées par clé,\
  \item pour chaque \texttt{asin}, on cumule la somme des notes
        (\texttt{sum\_ratings}) et le nombre de notes
        (\texttt{count\_ratings}),\
  \item au changement de clé, on calcule la moyenne
        \texttt{sum\_ratings / count\_ratings} et on l'affiche,\
  \item à la fin, on traite la dernière clé.
\end{itemize}

La sortie finale \texttt{task2\_output.txt} contient pour chaque
film~:\\
\texttt{asin\textbackslash taverage\_rating}.

\subsection{Exploration des résultats}

Comme pour la tâche 1, un échantillon est testé sur 1000 lignes
puis le résultat complet est enregistré dans \texttt{task2\_output.txt}.
Un tri décroissant sur la seconde colonne permet de voir les films
ayant les meilleures notes moyennes (en haut du classement).

\section{Tâche 3~: extraction de mots-clés fréquents}

\subsection{Objectif}

La troisième tâche ajoute une dimension \textbf{textuelle} en
analysant le champ \texttt{reviewText}. L'idée est d'extraire,
pour chaque film, les mots les plus fréquents dans les critiques,
après nettoyage et suppression des mots vides (\emph{stopwords}).

\subsection{Prétraitement du texte}

Le mapper de la tâche 3 (
\texttt{mapper\_task3.py}) utilise la bibliothèque NLTK pour~:
\begin{itemize}
  \item télécharger les ressources nécessaires~: stopwords anglais
        et tokeniseur (\texttt{punkt}),\
  \item définir une fonction \texttt{clean\_text(text)} qui~:
        \begin{itemize}
          \item met le texte en minuscules,\
          \item tokenise le texte en mots,\
          \item enlève les mots non alphabétiques,\
          \item enlève les mots vides (stopwords anglais).
        \end{itemize}
\end{itemize}

\subsection{Mapper de la tâche 3}

Pour chaque ligne JSON, le mapper~:
\begin{itemize}
  \item lit et parse la ligne en objet JSON,\
  \item récupère \texttt{asin} et \texttt{reviewText},\
  \item applique \texttt{clean\_text} au texte pour obtenir une
        liste de tokens pertinents,\
  \item pour chaque mot filtré, émet une paire de la forme~:
        \texttt{"asin:token\textbackslash t1"}.
\end{itemize}

Il s'agit donc d'un word count par \texttt{(asin, mot)}.

\subsection{Reducer de la tâche 3}

Le reducer \texttt{reduce\_task3.py} est un agrégateur classique~:
\begin{itemize}
  \item il lit des lignes \texttt{"asin:token\textbackslash tcount"},\
  \item les entrées étant triées, il maintient un \texttt{current\_key}
        (ici \texttt{"asin:token"}) et un compteur,\
  \item pour chaque changement de clé, il affiche la clé et la somme
        des occurrences,\
  \item à la fin, il affiche le dernier résultat.
\end{itemize}

Le fichier de sortie \texttt{task3\_output.txt} contient donc
pour chaque couple (film, mot) le nombre total d'occurrences du mot
dans les critiques de ce film.

\section{Tâche 4~: jointure des notes moyennes et des mots-clés}

\subsection{Objectif}

La quatrième tâche combine les résultats des tâches 2 et 3 pour
obtenir, pour chaque film~:
\begin{itemize}
  \item son identifiant \texttt{asin},\
  \item sa note moyenne (issue de \texttt{task2\_output.txt}),\
  \item la liste de ses mots-clés les plus fréquents (issue de
        \texttt{task3\_output.txt}).
\end{itemize}

On réalise pour cela une \textbf{reduce-side join}~: les deux fichiers
sont concaténés puis passés dans un même mapper qui les marque
respectivement comme données de notes (``R'') ou de mots-clés (``K'').

\subsection{Mapper de la tâche 4}

Le fichier \texttt{mapper\_task4.py} lit des lignes issues soit de
\texttt{task2\_output.txt}, soit de \texttt{task3\_output.txt}~:
\begin{itemize}
  \item si la partie clé (avant la tabulation) contient un deux-points
        (\texttt{:}), on la considère comme provenant de la tâche 3
        (format \texttt{asin:keyword}) et on émet~:\\
        \texttt{"asin\textbackslash tK:keyword:count"},\
  \item sinon, elle provient de la tâche 2 (format \texttt{asin\textbackslash trating})
        et on émet~:\\
        \texttt{"asin\textbackslash tR:rating"}.
\end{itemize}

Ainsi, toutes les données relatives à un même \texttt{asin} seront
réunies par le framework Hadoop sur un seul reducer.

\subsection{Reducer de la tâche 4}

Le reducer \texttt{reduce\_task4.py} reçoit, pour un \texttt{asin}
donné, une liste de valeurs préfixées par \texttt{R:} (une seule
note moyenne) ou \texttt{K:} (plusieurs couples mot--compte). Il~:
\begin{itemize}
  \item initialise \texttt{rating = None} et \texttt{keywords = []},\
  \item pour chaque valeur commencant par \texttt{R:}, enregistre
        la note moyenne,\
  \item pour chaque valeur \texttt{K:mot:compte}, ajoute \texttt{(mot, compte)}
        à la liste des mots-clés,\
  \item au changement d'\texttt{asin}, trie la liste des mots-clés
        par fréquence décroissante, construit une chaîne de la forme
        \texttt{"mot1(c1), mot2(c2), ..."} et affiche~:\\
        \texttt{"asin\textbackslash trating\textbackslash tmot1(c1), ..."},\
  \item répète l'opération pour la dernière clé.
\end{itemize}

Le résultat \texttt{task4\_output.txt} fournit ainsi une vision
synthétique par film~: pour chaque \texttt{asin}, on voit la note
moyenne et les mots dominants utilisés dans les critiques.

\section{Conclusion}

Ce mini-projet MapReduce sur les critiques Amazon Movies \\& TV
permet de mettre en pratique plusieurs compétences clés~:
\begin{itemize}
  \item installation et configuration d'un environnement Hadoop
        local (Java, Hadoop, variables d'environnement),\
  \item préparation d'un jeu de données volumineux (JSON compressé
        transformé en JSON Lines),\
  \item écriture de mappers et reducers en Python pour des tâches
        d'agrégation (counts, moyennes),\
  \item traitement de texte distribué pour l'extraction de mots-clés,\
  \item réalisation d'une \emph{reduce-side join} entre deux flux
        de résultats MapReduce.
\end{itemize}

Les sorties produites (nombre de critiques, note moyenne, mots-clés
les plus fréquents par film, jointure finale) constituent une base
précieuse pour comprendre la popularité des films, la satisfaction
des spectateurs et les thèmes récurrents évoqués dans leurs avis.

\end{document}
